\section[Introduction]{Introduction}

\begin{frame}{Introduction}
	Les Réseaux de Neurones sont une \textbf{famille} d'algorithmes d'apprentissage automatique flexibles permettant de s'adapter à des cas supervisés, non supervisés ou par renforcement.
	
	\vspace{0.3cm}
	
	Même si la technologie est ancienne, les réseaux de neurones ont eu un regain de popularité dans les années 2010 grâce:
	\begin{itemize}
		\item \`a l'émergence de la digitalisation des données et de l'informatique en ``Cloud'' qui permet d'accéder à de vastes quantités de données.
		\pause
		\item \`a l'explosion des moyens matériels de calculs, notamment les puces informatiques dédiées aux réseaux de neurones issues des cartes graphiques pour les jeux vidéo.
		\pause
		\item \`a l'accélération de la recherche scientifique dans le domaine de l'intelligence artificielle.
	\end{itemize}
	
\end{frame}
	
\begin{frame}{Contenu du cours}
	Au programme de ce cours:
	\begin{itemize}
		\pause
		\item Algorithme du Perceptron.
		\pause
		\item Résolution de problèmes non linéaires.
		\pause
		\item Algorithme de la rétro-propagation du gradient.
		\pause
		\item Fonctions de perte, optimisation et initialisation
	\end{itemize}
\end{frame}
